\section{Анализ предметной области и постановка задачи}

\subsection{Классы решений для контейнерной оркестрации}

В индустрии можно выделить несколько устойчивых классов средств.

\textbf{Кластерные оркестраторы} (прежде всего Kubernetes) ориентированы на крупные установки: декларативные манифесты, горизонтальное масштабирование, сетевые политики, хранилища и операторы. Недостатки для небольших проектов --- высокая кривая обучения, необходимость сопровождения control plane, иной цикл выпуска по сравнению с «одним compose-файлом на сервис''.

\textbf{Панели управления Docker} (в т.ч. Portainer) упрощают операции через UI и API поверх уже запущенного Docker. Они хорошо закрывают задачу визуализации и точечных действий, но реже выступают единой библиотекой для описания инфраструктуры из кода приложения или сборочного конвейера с тем же уровнем типобезопасности и переиспользования, что у доменной модели на языке разработки.

\textbf{Легковесный сценарий Docker Compose} сохраняет простоту и прозрачность: один файл описывает сервисы, сети и тома. Проблема возникает при нескольких хостах: нет встроенного центра реестра узлов, единого API публикации, согласованной политики бэкапа томов и управляемой публикации TCP-портов наружу без ручной правки прокси.

\subsection{Сравнительный анализ Chronos, Kubernetes и Portainer}

Таблица~\ref{tab:compare-chronos-k8s-portainer} обобщает отличия по наиболее значимым для данной работы критериям. Оценки носят качественный характер: конкретные редакции Portainer и дистрибутивов Kubernetes могут расширять отдельные пункты плагинами.

{\small
\begin{longtable}{@{}p{0.22\linewidth}p{0.25\linewidth}p{0.25\linewidth}p{0.25\linewidth}@{}}
\caption{Сравнение Chronos, Kubernetes и Portainer} \label{tab:compare-chronos-k8s-portainer} \\
\toprule
\textbf{Критерий} & \textbf{Chronos} & \textbf{Kubernetes} & \textbf{Portainer} \\
\midrule
\endfirsthead
\multicolumn{4}{c}{\tablename\ \thetable{} --- продолжение} \\
\toprule
\textbf{Критерий} & \textbf{Chronos} & \textbf{Kubernetes} & \textbf{Portainer} \\
\midrule
\endhead
\midrule
\multicolumn{4}{r}{\small\textit{Продолжение на следующей странице}} \\
\endfoot
\endlastfoot

Целевая аудитория &
Команды и индивидуальные разработчики, уже использующие Compose; небольшие стенды. &
Организации с выделенной платформенной командой и кластерами от нескольких узлов. &
Администраторы и разработчики, управляющие Docker через UI/API без обязательной кластеризации. \\
\addlinespace

Модель описания сервисов &
Docker Compose (YAML) и/или Fluent API на C\# в библиотеке Core. &
Pod/Deployment/\newline Service/Ingress и др.; часто Helm/Kustomize. &
В основном UI и Docker Compose-стеки; «инфраструктура как код'' чаще внешняя по отношению к Portainer. \\
\addlinespace

Несколько хостов без K8s &
Ядро: Master + набор Agent; единый UI и политики. &
Требуется полноценный кластер Kubernetes. &
Есть Portainer Agent на узлах; фокус на обзоре/действиях, не на единой модели публикации compose из .NET. \\
\addlinespace

Программируемость из кода приложения &
Прямое использование Chronos.Core в сборках, тестах, консольных утилитах. &
Клиенты API (\texttt{kubectl}, client-go, CDK8s и т.д.); иная модель абстракций. &
REST/API Portainer; нативной Fluent-модели compose под .NET нет. \\
\addlinespace

Политики снимков томов в объектное хранилище &
Встроены: политики на Master, снимок на Agent, S3-совместимый бакет, ретенция \texttt{MaxCopies}. &
VolumeSnapshot API, внешние операторы бэкапа; настройка сложнее. &
Есть сценарии резервного копирования томов (зависит от редакции); нет той же связки «политика Master $\rightarrow$ агент'' из коробки. \\
\addlinespace

Динамический TCP edge &
Реестр маршрутов Master, генерация HAProxy, диапазон listen-портов. &
Service type LoadBalancer / Ingress-контроллеры; своя модель. &
Не является центральной функцией продукта; настраивается внешними средствами. \\
\addlinespace

Самописные проверки и фоновые задачи в рантайме compose &
Кодовые тесты (\texttt{[Test]}) и job-ы (\texttt{[Job]}) в подключаемых сборках; запуск на агенте после старта и по расписанию. &
Probes (liveness/readiness), CronJob как отдельная абстракция, не «внутри'' compose-сервиса в смысле Chronos. &
Healthcheck в compose; произвольный C\# на агенте в модели Chronos не предусмотрен. \\
\addlinespace

Порог входа и сопровождение &
Ниже, чем у production Kubernetes; нужны Docker и понимание compose. &
Высокий; требуется экспертиза по сети, хранилищам, обновлениям control plane. &
Средний для UI-операций; кластерные сценарии усложняются. \\
\addlinespace

Отказоустойчивость control plane (типичный OSS-сценарий) &
MVP: один Master; задел под лидерство. &
Распределённый etcd, несколько control plane (в зависимости от установки). &
Зависит от установки Portainer Server; обычно репликация/reverse proxy настраиваются отдельно. \\
\bottomrule

\end{longtable}
}

\subsection{Обзор источников и нормативной базы}

Для обоснования архитектурных решений и выбора стека опирались на открытые спецификации, документацию вендоров и признанные монографии по контейнеризации и платформе .NET~\cite{docker_get_started,docker_engine_docs,matthias2018docker,boettiger2015docker,dotnet_docs,aspnetcore_docs,smith2022modern_web,griffiths2021csharp,freeman2023proaspnet}.

\textbf{Docker и Compose.} Обзор контейнерной модели и жизненного цикла образов приведён в документации Docker Engine~\cite{docker_engine_docs,docker_get_started} и в практическом руководстве~\cite{matthias2018docker}; обзорный характер контейнеризации для воспроизводимых сред рассмотрен в~\cite{boettiger2015docker}. Формат манифеста и семантика полей \texttt{services}, \texttt{networks}, \texttt{volumes}, \texttt{depends\_on} зафиксированы в официальной спецификации Compose~\cite{docker_compose_spec}; Chronos сознательно не дублирует планировщик контейнеров, а генерирует совместимый YAML и вызывает штатный \texttt{docker compose} на агенте.

\textbf{Платформа .NET.} Серверные компоненты Chronos (Master, Agent) реализованы на ASP.NET Core; ориентиры по архитектуре веб-приложений, Minimal API и персистентности данных --- в документации Microsoft~\cite{dotnet_docs,aspnetcore_docs,aspnetcore_minimal_apis,ef_core_docs} и в учебных изданиях по C\# и ASP.NET Core~\cite{griffiths2021csharp,freeman2023proaspnet,smith2022modern_web}.

\textbf{Kubernetes.} Руководства по Deployment, Service, Ingress и VolumeSnapshot~\cite{burns2018kubernetes} дают ориентиры по тому, какие задачи в крупных кластерах решаются декларативно; сопоставление с Chronos выполняется в таблице~\ref{tab:compare-chronos-k8s-portainer}. Для ВКР важно зафиксировать: Kubernetes остаётся эталоном по полноте модели, но избыточен для сценария «несколько VPS с Docker'' без выделенной платформенной команды.

\textbf{HAProxy.} Режим \texttt{mode tcp}, директивы \texttt{bind} и \texttt{server}, а также ограничения парсера конфигурации (кодировка, переносы строк) описаны в официальной документации HAProxy~\cite{haproxy_docs}. Динамическое подключение фрагмента конфигурации через \texttt{include} и сигнал перезагрузки --- распространённый приём; в работе он применён для реестра маршрутов Chronos.

\textbf{Portainer и смежные панели.} Документация API и модели агентов Portainer иллюстрирует подход «центральный сервер + агенты на хостах''; Chronos перенимает идею регистрации узлов, но добавляет доменную модель compose в .NET и политики бэкапа томов, не претендуя на универсальную администрацию всего Docker-хоста.

\textbf{Наблюдаемость.} Модель логов Loki (лейблы, потоки, запросы LogQL) и роль Grafana как визуального слоя описаны в документации Grafana Labs. Выбор push-логов по HTTP упрощает развёртывание на малых стендах по сравнению с полноценным стеком на базе Prometheus для каждого узла.

\subsection{Позиционирование Chronos}

Chronos занимает нишу между «только Compose на каждой машине'' и полноценным Kubernetes: сохраняется формат и инструментарий Compose, добавляется программируемое описание через .NET-библиотеку, центральный узел и агенты на хостах с Docker.

Вдохновляющим примером для постановки задачи послужили облачные практики (в духе Azure Resource Manager / Bicep и родственных подходов), где параметры среды задаются декларативно и проверяются до развёртывания. В Chronos аналогичную роль выполняет связка Fluent API и валидатора с последующей генерацией YAML и выкладкой на агент.

\subsection{Проблемы исходного сценария Compose-only}

Без управляющего слоя типичны следующие трудности:
\begin{itemize}[leftmargin=1.5cm]
  \item ручной или нестандартизованный деплой на разных хостах;
  \item отсутствие единого API и журнала операций;
  \item сложность маршрутизации внешнего TCP-трафика к сервисам на разных машинах;
  \item отсутствие единой политики снимков томов и хранения бэкапов в объектном хранилище.
\end{itemize}

\subsection{Функциональные требования}

\begin{enumerate}[leftmargin=1.5cm]
  \item Описание compose-проекта из кода (Fluent API) и/или YAML с валидацией.
  \item Публикация проекта на удалённый агент по HTTP, управление жизненным циклом (start/stop/restart).
  \item Регистрация агентов на Master, heartbeat, отображение состояния в UI.
  \item Декларативные проверки и задания на стороне агента (после старта и по расписанию).
  \item Резервное копирование томов по политикам Master с выгрузкой в S3-совместимое хранилище.
  \item Динамическое добавление и удаление TCP-маршрутов с генерацией конфигурации HAProxy.
  \item Централизованный сбор логов приложений в Loki и визуализация в Grafana.
\end{enumerate}

\subsection{Нефункциональные требования}

\begin{itemize}[leftmargin=1.5cm]
  \item воспроизводимость поведения в локальной и контейнерной среде;
  \item расширяемость (новые сценарии без переписывания ядра);
  \item наблюдаемость (логи, метрики, трассировка запросов к API);
  \item устойчивость к рестартам в модели single-master MVP;
  \item контроль секретов и редактирование чувствительных фрагментов в логах.
\end{itemize}

\subsection{Выбор технологического стека}

\begin{itemize}[leftmargin=1.5cm]
  \item \textbf{.NET 8} --- единая платформа для Core, Agent и Master;
  \item \textbf{Entity Framework Core + PostgreSQL} --- метаданные кластера, аудит, политики бэкапа;
  \item \textbf{React + Vite + TypeScript} --- веб-интерфейс;
  \item \textbf{HAProxy} --- TCP-прокси с динамически подключаемым фрагментом конфигурации;
  \item \textbf{Loki + Grafana} --- операционные логи;
  \item \textbf{S3 API (MinIO и др.)} --- хранение снимков томов.
\end{itemize}

\subsection{Краткая схема модулей (для последующих глав)}

Для дальнейшего изложения зафиксируем декомпозицию:
\begin{enumerate}[leftmargin=1.5cm]
  \item \textbf{Chronos.Core} --- доменная модель compose, генерация и разбор YAML, клиент агента, локальный тестер, артефакты деплоя, исполнение кодовых jobs и тестов.
  \item \textbf{Chronos.Agent} --- исполнение \texttt{docker compose}, хранение проектов, снимки томов, планировщик проверок, связь с Master.
  \item \textbf{Chronos.Master} --- персистентность, UI API, прокси к агентам, HAProxy-реестр, фоновая репликация томов, раздача собранного SPA.
\end{enumerate}

\subsection{Постановка задачи}

Требуется спроектировать и реализовать программный комплекс, который принимает декларативное описание контейнерных сервисов (из кода или YAML), публикует его на удалённых агентных узлах, обеспечивает централизованное управление запуском, наблюдаемость, политики резервного копирования томов и проксирование внешнего TCP-трафика к сервисам на хостах агентов через HAProxy, сохраняя совместимость с Docker Compose.

\subsection{Цели, показатели качества и границы MVP}

\textbf{Цель разработки} сформулирована во введении: получить воспроизводимый контур управления Compose-проектами на нескольких узлах с веб-интерфейсом и минимально достаточной наблюдаемостью.

\textbf{Показатели качества} (в качественной форме, пригодной для проверки на стенде):
\begin{itemize}[leftmargin=1.5cm]
  \item \emph{Корректность API:} ответы согласованы с фактическим состоянием агента и БД; ошибки валидации возвращаются до побочных эффектов (например, до записи конфликтующего TCP-маршрута).
  \item \emph{Надёжность сценария TCP:} после обновления \texttt{chronos-tcp.cfg} и перезагрузки HAProxy внешнее TCP-соединение доходит до опубликованного на хосте агента порта в типовой конфигурации Docker Desktop / Linux.
  \item \emph{Воспроизводимость бэкапа:} политика на Master инициирует снимок не чаще заданного интервала; число объектов в бакете не превышает \texttt{MaxCopies} после процедуры prune.
  \item \emph{Наблюдаемость:} критические ошибки API и фоновых служб видны в Loki и позволяют локализовать узел (Master или Agent) по лейблам.
\end{itemize}

\textbf{Границы MVP} (что сознательно не входит в первую версию, но оставляет направление развития):
\begin{itemize}[leftmargin=1.5cm]
  \item отказоустойчивый кластер Master и автоматический failover;
  \item полноценный RBAC и интеграция с корпоративными IdP;
  \item автоматическое обнаружение публикуемых портов compose без участия оператора при создании TCP-маршрута;
  \item формальные нагрузочные тесты с отчётом SLA (рекомендуются как отдельный этап после защиты прототипа).
\end{itemize}
